\FloatBarrier{}
\section{Introducción}%
\label{sec:introduction}

En la actualidad la organización eficiente de la información vinculada a viajes se ha convertido en un desafío.
Es común que para el armado de un viaje se recurran a varios proveedores distintos de servicios turísticos.
Si se realiza una compra en el sitio web de la aerolínea (por ejemplo a través de Lufthansa) el cliente recibirá en su
bandeja de entrada un correo electrónico con todos los datos de su reserva.
Cuando luego se realice la reserva del hotel es probable que utilice un sitio web diferente (por ejemplo a través de
Booking.com) y obtenga otro correo electrónico distinto con los datos de la estadía.
También es posible que durante el viaje se programen vuelos locales o regionales que involucren otras aerolíneas
distintas (por ejemplo Ryanair).

Para poder consultar información importante como el aeropuerto de origen del vuelo o el horario de check-out del hotel
el usuario tiene que estar consultando correos electrónicos distintos.
Aunque estas empresas proveedoras de servicios turísticos ofrecen también aplicaciones móviles, tampoco se consigue una
vista integral ya que el usuario deberá descargar y consultar una aplicación para el vuelo de ida, otra aplicación
diferente para el hotel y quizás otra aplicación distinta para algún vuelo regional.

Además, al momento de darle un vistazo a los correos electrónicos o las distintas aplicaciones, cada una presenta una
interfaz distinta donde es probable que el aeropuerto de origen del vuelo en el mail de una aerolínea esté arriba a la
izquierda mientras que en otra aparezca en el centro.
Estructuras distintas para diferentes proveedores obligan a que, para lograr una vista integral de la información, el
usuario termine copiando esa información en otra parte, corriendo el riesgo de cometer un error o de contar con
información desactualizada.

Es por eso que es importante contar con una aplicación capaz de extraer información relevante de correos electrónicos de
reservas independientemente de la estructura o el remitente del correo para luego proceder a la creación de un
itinerario del viaje integral.

El presente proyecto se centrará en la extracción de información relevante contenida exclusivamente en correos
electrónicos de reservas de tickets aéreos y de estadías de hotel.
Si bien se excluyen otros tipos de reserva (como alquileres de vehículos, pasajes de tren o reservas de restaurantes)
se buscará diseñar una solución extensible en caso de que se desee ampliar el dominio en el futuro.

Asimismo, el alcance del trabajo abarca el desarrollo de un prototipo funcional que permita validar las técnicas
aplicadas de extracción de información utilizando una solución web simple que simule ser la vista del usuario de la
aplicación de gestión de itinerarios.
No se contemplarán requisitos funcionales y no funcionales propios de una aplicación real orientada a un gran volumen
de usuarios, tales como usabilidad, alta disponibilidad y seguridad.